% TODO checklist:
% - [x] Inspected src/workers/jobs_worker.py (heartbeat, signal handling)
% - [x] Inspected db/redis_cache/job_store.py (atomic operations, orphan cleanup)
% - [x] Inspected src/background.py (scheduler resilience)
% - [x] Analyzed Q&A.md Q41 for operational strengths/weaknesses
% - [x] Reviewed health probe integration patterns

\section{High Availability and Fault-Tolerance Patterns}
\label{sec:ha-ft}

Enterprise deployments require the job system to operate reliably under adverse conditions: worker crashes, network partitions, database unavailability, and infrastructure restarts. This section describes the fault-tolerance patterns implemented in the Cineca Agentic Platform's job and worker subsystems.

\subsection{Worker Heartbeat and Dead-Worker Detection}

Workers maintain liveness signals through periodic heartbeat updates to the job's \texttt{updated\_at} timestamp:

\begin{lstlisting}[language=Python, caption={Heartbeat mechanism for liveness detection.}]
async def _heartbeat_loop(self, job_uuid: UUID, jobs_service: JobsService):
    """Periodically update job timestamp to indicate worker is alive."""
    while True:
        await asyncio.sleep(self.heartbeat_interval)
        try:
            jobs_service.repo.touch_job(job_uuid)
            logger.debug(f"Heartbeat for job {job_uuid}")
        except Exception as e:
            logger.warning(f"Heartbeat failed: {e}")
\end{lstlisting}

Dead workers can be detected by monitoring jobs in RUNNING status whose \texttt{updated\_at} exceeds a threshold (e.g., 3× heartbeat interval). A monitoring process or background task can identify such jobs for recovery.

\subsection{Orphaned Job Recovery}

Jobs left in RUNNING status after a worker crash require recovery. The recovery process:

\begin{enumerate}
    \item \textbf{Detection}: Query for RUNNING jobs with stale heartbeats.
    \item \textbf{Requeue}: Transition job back to QUEUED status.
    \item \textbf{Execution}: Normal worker picks up the requeued job.
\end{enumerate}

\begin{lstlisting}[language=Python, caption={Orphaned job recovery (conceptual).}]
async def recover_orphaned_jobs(threshold_seconds: int = 60):
    """Recover jobs stuck in RUNNING with stale heartbeats."""
    cutoff = datetime.utcnow() - timedelta(seconds=threshold_seconds)
    
    orphaned = await jobs_service.find_stale_running_jobs(cutoff)
    
    for job in orphaned:
        logger.warning(f"Recovering orphaned job {job.id}")
        # Transition back to queued for retry
        await jobs_service.transition_status(
            job.id, from_status="running", to_status="queued"
        )
        # Re-enqueue to Redis
        await jobs_cache.queue_push_job(job.type, job.id)
\end{lstlisting}

\subsection{Retry and Backoff Strategies}

The current implementation does not include automatic retry for failed jobs. Table~\ref{tab:retry-recommendations} outlines recommended patterns for production deployments.

\begin{table}[ht]
\centering
\caption{Recommended retry patterns for job failures.}
\label{tab:retry-recommendations}
\small
\begin{tabular}{lll}
\hline
\textbf{Failure Type} & \textbf{Strategy} & \textbf{Configuration} \\
\hline
Transient errors & Exponential backoff & Base: 1s, Max: 60s, Factor: 2 \\
Resource exhaustion & Fixed delay retry & 30s delay, max 3 retries \\
Permanent failures & No retry & Move to dead-letter queue \\
Cancellation & No retry & Final state, no recovery \\
\hline
\end{tabular}
\end{table}

\subsection{Graceful Degradation}

The system implements graceful degradation when dependencies become unavailable:

\subsubsection{Redis Unavailability}

When Redis is unavailable, the job system:
\begin{itemize}
    \item Falls back to in-memory storage (development/testing mode)
    \item Returns 503 Service Unavailable for job operations
    \item Health probes report DEGRADED status
\end{itemize}

\subsubsection{PostgreSQL Unavailability}

When PostgreSQL is unavailable:
\begin{itemize}
    \item Workers pause job processing
    \item In-flight jobs continue but cannot persist results
    \item Health probes report UNHEALTHY status
    \item API endpoints return 503 for database-dependent operations
\end{itemize}

\subsection{Health Probe Integration}

The job system integrates with Kubernetes health probes for automated recovery:

\begin{lstlisting}[language=Python, caption={Health probe integration for job system.}]
async def check_job_system_health() -> ComponentCheck:
    """Check job system health for readiness probe."""
    checks = {}
    
    # Check Redis connectivity
    try:
        redis = await get_async_redis()
        await redis.ping()
        checks["redis"] = "ok"
    except Exception as e:
        checks["redis"] = f"error: {e}"
        return ComponentCheck(ok=False, status="error", details=checks)
    
    # Check queue depths
    for job_type in allowed_types:
        depth = await redis.llen(f"jobs:queue:{job_type}")
        checks[f"queue_{job_type}"] = depth
    
    # Check for stuck jobs (optional)
    stuck_count = await count_stale_running_jobs()
    if stuck_count > 0:
        checks["stuck_jobs"] = stuck_count
        return ComponentCheck(ok=True, status="degraded", details=checks)
    
    return ComponentCheck(ok=True, status="ok", details=checks)
\end{lstlisting}

\subsection{Index Consistency Maintenance}

Redis ZSET indexes can become inconsistent when job HASH keys expire but index entries remain. The orphan cleanup task (Section~\ref{subsec:scheduled-task-categories}) maintains index consistency by periodically scanning indexes and removing entries for non-existent jobs.

\subsection{Transaction Boundaries}

Critical operations use appropriate transaction boundaries:

\begin{itemize}
    \item \textbf{Job Creation}: Atomic creation of HASH + ZSET entries via Redis pipeline.
    \item \textbf{Status Updates}: Lua scripts for atomic compare-and-swap operations.
    \item \textbf{PostgreSQL Updates}: SQLAlchemy session transactions with rollback on failure.
\end{itemize}

\subsection{Operational Recommendations}

Based on the implemented patterns, Table~\ref{tab:ha-recommendations} provides operational recommendations for high-availability deployments.

\begin{table}[ht]
\centering
\caption{High availability operational recommendations.}
\label{tab:ha-recommendations}
\small
\begin{tabular}{lp{7cm}}
\hline
\textbf{Component} & \textbf{Recommendation} \\
\hline
Workers & Run 2+ worker replicas for redundancy \\
Redis & Use Redis Sentinel or Cluster for HA \\
PostgreSQL & Use managed service with automatic failover \\
Monitoring & Alert on queue backlog $>$ threshold \\
Recovery & Schedule orphan job recovery every 5 minutes \\
Backups & Enable daily backups with 7-day retention \\
\hline
\end{tabular}
\end{table}

\subsection{Known Limitations}

The current implementation has several limitations that should be addressed for mission-critical deployments:

\begin{enumerate}
    \item \textbf{No Automatic Retry}: Failed jobs remain in FAILED status; manual intervention or custom recovery is required.
    
    \item \textbf{Single Worker Assumption}: No distributed locking for multi-worker deployments; duplicate processing possible without Redis SETNX guards.
    
    \item \textbf{No Dead-Letter Queue}: Permanently failed jobs are not segregated for manual review.
    
    \item \textbf{Basic Priority Support}: Priority field exists but is not fully utilized in queue polling logic.
\end{enumerate}

Future enhancements should address these limitations through:
\begin{itemize}
    \item Configurable retry policies with exponential backoff
    \item Distributed locking via Redis SETNX or Redlock algorithm
    \item Dead-letter queue routing for failed jobs
    \item Priority-weighted queue polling
\end{itemize}

